

% -------------- Resume ---------------
\documentclass{resume}

\begin{document}

% --------- Contact Information -----------
\introduction[
    fullname={Phùng Quang Hưng},
    email={phunghung146@gmail.com},
    phone={0963624643},
    linkedin={linkedin.com/in/hunggpq}
    % github={github.com/sparkysundevil}
]
\summary{DevOps Engineer with hands-on experience in Kubernetes and AWS infrastructure. Focused on infrastructure automation, reliable application delivery, and continuous improvement.
}


% --------- Education ---------
\begin{educationSection}{Education}
    \educationItem[
        university={University of Information Technology - VNUHCM },
        graduation={2022 -- 2026},
        grade={8.75/10.0 CGPA},
        program={B.Sc., Computer Networks and Data Communications},
        coursework={Cloud Computing, DevOps, Network and System Administration, Distributed Systems},
    ]
\end{educationSection}
% --------- Work Experience ----------
\begin{experienceSection}{Work Experience}
    \experienceItem[
        company={FPT Software},
        position={DevOps Engineer},
        location={Ho Chi Minh City},
        duration={Nov 2025 -- Present}
    ]
    \begin{itemize}
        \vspace{-0.5em}
        \itemsep -6pt {}
        \item Contributed to migrating 20 services from Nomad on AWS EC2 to Amazon EKS, converting and standardizing YAML deployment manifests to platform standards.
        \item Set up AI-assisted skills for Nomad-to-Kubernetes YAML conversion and deployment, and configured GitHub Actions CI/CD using shared modules.
        \item Monitored production on EKS with CloudWatch and Datadog, troubleshooting incidents to keep services stable.
    \end{itemize}
    \experienceItem[
        company={Axon Active},
        position={Junior ICT Engineer},
        location={Ho Chi Minh City},
        duration={Aug 2025 -- Oct 2025}
    ]
    \begin{itemize}
        \vspace{-0.5em}
        \itemsep -6pt {}
        \item Provided IT support (L1/L2), account management, and endpoint troubleshooting.
    \end{itemize}

\end{experienceSection}
% --------- Projects -----------
\begin{experienceSection}{Academic Projects}
	\projectItem[
	title={ROOT CAUSE ANALYSIS IN MICROSERVICES},
	duration={June 2026},
	team={2},
	keyHighlight={Built MyRCA, an unsupervised two-step RCA framework for microservices via distributed tracing. \href{https://drive.google.com/drive/u/3/folders/1t9heWwOSkCfPsiXrWNROt5kCzmniqRRA}{\textcolor{blue}{\textit{Project Report}}}}
	]
	\begin{itemize}
		\vspace{-0.5em}
		\itemsep -6pt {}
		\item Built and deployed a microservices root-cause analysis framework using distributed tracing data and Kubernetes, enabling automated diagnosis of service failures.
		\item Implemented the anomaly detection pipeline in Python and ranked root-cause candidates using Causal \& Jaccard scoring, achieving F1 0.877 and A@1 0.909.
	\end{itemize}
	
	\projectItem[
	title={End-to-end DevOps Pipeline - Ranked 1st in class},
	duration={June 2025},
	team={2},
	keyHighlight={Deployed chat microservices app on AWS with Terraform, Jenkins \& ArgoCD. \href{https://github.com/MeetingTeam}{\textcolor{blue}{\textit{GitHub}}}}
	]
	\begin{itemize}
		\vspace{-0.5em}
		\itemsep -6pt {}
		\item Provisioned AWS infrastructure with Terraform and deployed a microservices application using a full end-to-end DevOps workflow.
		\item Configured CI/CD pipelines with Jenkins and ArgoCD and managed Kubernetes manifests to support automated delivery and environment consistency.
		
	\end{itemize}
	
	
	% \projectItem[
	%     title={Enterprise Network Infrastructure Design for O-UIT Company},
	%     duration={Fall 2024},
	%     team={4},
	%     keyHighlight=Design three-tier architect enterprise network.
	%     % Optional: Add report link if available and professional
	%     \href{https://docs.google.com/document/d/1ij7DAoHRcJV_3vilYMOIAESrTe84WP7t/edit?usp=sharing&ouid=110580317238341657252&rtpof=true&sd=true}{\textcolor{blue}{\textit{Project Report}}}
	% ]
	
\end{experienceSection}


% --------- Skills -----------
\begin{skillsSection}{Technical Skills}
    \skillItem[
        category={Cloud \& IaC},
        skills={AWS, Terraform}
    ] \\
    \skillItem[
        category={CI/CD \& Containers},
        skills={Git, GitHub, Jenkins, ArgoCD, Docker, Kubernetes}
    ] \\
    \skillItem[
        category={Cloud Platforms \& Monitoring},
        skills={AWS EC2, Amazon EKS, Nomad, CloudWatch, Datadog}
    ] \\
    \skillItem[
        category={Networking \& OS},
        skills={TCP/IP, VLAN, Routing, Linux}
    ] 
\end{skillsSection}







% --------- Awards -----------
\begin{experienceSection}{Awards and Certificates}
    \projectItemm[
        title={NET Challenge 2024 – Champion \& Team Leader},
        duration={Dec 2024},
        keyHighlight={Led the team to 1st Prize among 109 teams nationwide (Networking \& Cloud).}
    ] \\
    \projectItemm[
        title={Academic Scholarship – UIT-VNUHCM},
        duration={2023–2025},
        keyHighlight={Awarded for consecutive excellent academic performance.}
    ] \\
    \awardItem[
        title={AWS Certified Solutions Architect – Associate},
        duration={Jan 2026}
    ] \\
    \awardItem[
        title={AWS Certified AI Practitioner},
        duration={Sep 2026}
    ] \\
    \awardItem[
        title={TOEIC LR 810},
        duration={Sep 2025}
    ] \\
    \awardItem[
        title={CCNA},
        duration={Aug 2024}
    ] 
\end{experienceSection}

\end{document}
